% Copyright Ben Paul Wise. All Rights Reserved.
\documentclass[11pt,letterpaper]{article}
\usepackage[letterpaper, margin=1in]{geometry}
\usepackage{enumitem}
\usepackage{booktabs}
\usepackage{parskip}
\usepackage{titlesec}
\usepackage{microtype}
\usepackage{hyperref}
\usepackage{xcolor}

\hypersetup{colorlinks=true, linkcolor=blue!60!black}

\titleformat{\section}{\large\bfseries}{}{0em}{}[\titlerule]
\titleformat{\subsection}{\normalsize\bfseries}{}{0em}{}

\newcommand{\file}[1]{\texttt{#1}}
\newcommand{\code}[1]{\texttt{#1}}
\newcommand{\bug}{\textcolor{red}{\textbf{BUG}}}
\newcommand{\pbug}{\textcolor{orange}{\textbf{POTENTIAL BUG}}}
\newcommand{\ineff}{\textcolor{blue}{\textbf{INEFFICIENCY}}}
\newcommand{\style}{\textcolor{gray}{\textbf{STYLE}}}

\title{\textbf{IrrGo C++ Code Review}}
\author{Ben Paul Wise}
\date{\today}

\begin{document}
\maketitle
\tableofcontents
\newpage

% ────────────────────────────────────────────────────────────────────────────
\section{Group 1 — Definite Bugs}
% ────────────────────────────────────────────────────────────────────────────

\subsection{\bug\ \ \file{src/Searcher.h}:26 — \code{terminalCount} is not thread-safe}

\code{terminalCount} is a \code{static} member variable written by the search
thread and reset by the launching thread with no synchronisation.  If two
searches run concurrently (which the current guard prevents, but is fragile),
the counter is silently corrupted.  Even with the guard, the reset in the
main thread and the increment in the worker thread are a data race under the
C++ memory model.  Fix: use \code{std::atomic<int>}.

\subsection{\bug\ \ \file{lib/Game.cpp} — Pass move leaves \code{col} default-initialised to 0}

The \code{Move} struct has five fields: \code{turn}, \code{color},
\code{nodeId}, \code{row}, \code{col}.  The pass-move push uses only four
initialisers, so \code{col} is default-initialised to \code{0} instead of
\code{-1}.  Any code that reads \code{col} from a pass move and compares it
against \code{-1} will behave incorrectly.

\subsection{\bug\ \ \file{gui/MainWindow.cpp} — Detached \code{std::thread} holds raw \code{this}}

All four search launchers (\code{onSuggestGo}, \code{onSuggestMctsGo},
\code{onPlayNegamaxGo}, \code{onPlayMctsGo}) detach a \code{std::thread}
whose lambda captures \code{this} by value.  If \code{MainWindow} is
destroyed before the thread completes, all member-variable accesses in the
lambda are undefined behaviour.  The generation-counter cancellation
mechanism reduces the risk but does not eliminate it.  Fix: use a
\code{std::shared\_ptr} sentinel or join on destruction.

% ────────────────────────────────────────────────────────────────────────────
\section{Group 2 — Potential Bugs}
% ────────────────────────────────────────────────────────────────────────────

\subsection{\pbug\ \ \file{lib/Graph.h}:15 — \code{node()} has no bounds check}

\code{Graph::node(int id)} indexes \code{nodes\_[id]} without validating
that \code{id} is in range.  Any caller passing an invalid ID causes
undefined behaviour.  An \code{assert} or a bounds check with a clear error
message would make failures diagnosable.

\subsection{\pbug\ \ \file{lib/Graph.cpp} — \code{addEdge()} does not guard against duplicates}

\code{addEdge(a, b)} appends \code{b} to \code{a.neighbors} and vice versa
unconditionally.  If called twice for the same pair, the neighbor lists
contain duplicates.  Every algorithm that iterates neighbors (DVR BFS, liberty
counting, edge export in \file{MainWindow\_save.cpp}) will then process each
edge twice.

\subsection{\pbug\ \ \file{src/mcts.cpp}:45--54 — UCT score called with zero visit count}

\code{uctScore()} divides by \code{child.visitCount} without checking for
zero.  While the current \code{treePolicy()} only calls it after all children
are expanded (each having been through at least one \code{backup()}), the
function itself provides no guard.  A future refactor or a subtle ordering
change would cause a division-by-zero crash.  Add an assertion or early
return.

\subsection{\pbug\ \ \file{gui/BoardWidget.cpp}:153--167 — \code{nodeAt()} returns stale index}

The nearest-node search returns the raw index \code{found} from the game's
node list.  If the board has been regenerated between the mouse event and the
lookup (a race that the current single-threaded UI makes unlikely but not
impossible), \code{found} may exceed the new node count.  The result is
passed to \code{game\_->placeStone()} without a bounds check.

\subsection{\pbug\ \ \file{gui/BoardWidget.cpp}:293--310 — Manhattan distance meaningless for irregular graphs}

The neighbourhood-highlight code computes
$|r_a - r_b| + |c_a - c_b|$ to decide whether two nodes are ``nearby.''
For a \code{RectangularGraph} this approximates physical distance, but for an
\code{IrregularGraph} row/col are base-grid coordinates and bear no
consistent relationship to the actual layout.  The displayed neighbourhood
radius will be visually wrong for irregular boards.

\subsection{\pbug\ \ \file{gui/MainWindow\_save.cpp}:26--132 — XML write errors are silently dropped}

After the initial \code{file.open()} check, no error is ever inspected.
\code{QXmlStreamWriter} and \code{QFile} can fail mid-write (disk full,
permissions, etc.) without the caller knowing.  The saved file would be
truncated or malformed with no warning to the user.  Check
\code{xml.hasError()} and \code{file.error()} before closing.

\subsection{\pbug\ \ \file{gui/MainWindow\_save.cpp}:34 — Open-coded graph type check will silently mis-classify new graph types}

\code{isRect} is set by a \code{dynamic\_cast} to \code{RectangularGraph}.
If a third graph type is ever added, it will silently be saved as
\code{rectangularP val="False"} (i.e.\ irregular), which may be incorrect.
A virtual method such as \code{Graph::isRectangular()} would be more robust.

% ────────────────────────────────────────────────────────────────────────────
\section{Group 3 — Inefficiencies}
% ────────────────────────────────────────────────────────────────────────────

\subsection{\ineff\ \ \file{lib/Game.cpp} — \code{getLegalMoves()} is O($N^2$)}

\code{isLegalPlacement()} copies the entire board into \code{lp\_board\_}
on every call (O($N$)).  \code{getLegalMoves()} calls it once per node, making
the total O($N^2$).  During MCTS rollouts this is the innermost hot loop.
A diff-based or incremental legality check would give a large speedup,
especially on 19$\times$19 or 21$\times$21 boards.

\subsection{\ineff\ \ \file{src/mcts.cpp}:94 — Rollout clones the full game state}

Each MCTS rollout calls \code{node.game->clone()}, duplicating the entire
board for every simulation.  For a game with $N$ nodes this allocates
and copies $O(N)$ data per rollout, millions of times per search.  An
in-place playout with undo (or a lighter-weight game representation) would
reduce allocator pressure significantly.

\subsection{\ineff\ \ \file{gui/BoardWidget.cpp}:155--157 — \code{nodeAt()} is O($N$) per mouse event}

Every hover and click linearly scans all nodes to find the nearest one.
On a 21$\times$21 irregular board (441 nodes) this is fast enough, but
for larger or denser graphs it becomes perceptible.  A spatial grid or
k-d tree built once after board generation would reduce this to O(1).

\subsection{\ineff\ \ \file{gui/BoardWidget.cpp}:293--310 — Neighbourhood highlight is O($N^2$) per paint}

For each node being highlighted, all other nodes are scanned to compute
a distance.  This runs on every \code{paintEvent()}, including during
animations.  The neighbourhood size for each node should be computed once
after board generation and cached.

\subsection{\ineff\ \ \file{lib/IrregularGraph.cpp} — Edge validation is O($E \cdot N$) per candidate}

\code{canAddEdge()} checks crossing and proximity against all existing edges
and all nodes.  During graph construction this is called repeatedly in a
greedy loop, making total construction O($E^2 N$) in the worst case.  A
spatial index on node positions would reduce crossing tests to the relevant
neighbourhood.

\subsection{\ineff\ \ \file{lib/DVR.cpp} — Two independent BFS passes could be merged}

\code{DVR} runs two separate multi-source BFS passes (\code{distMine} and
\code{distOpp}).  A single interleaved BFS that assigns ownership to each
node on first discovery would halve the traversal work and improve cache
locality.

% ────────────────────────────────────────────────────────────────────────────
\section{Group 4 — Style and Maintainability}
% ────────────────────────────────────────────────────────────────────────────

\subsection{\style\ \ \file{src/mcts.cpp}:174--175, 196--197 — Debug prints left in production code}

\code{fprintf(stderr, ...)} and \code{printf(...)} are called at the end of
both MCTS overloads on every search.  These should be removed or placed
behind a compile-time flag before any non-debug build.

\subsection{\style\ \ \file{src/AbsGame.cpp} — Empty translation unit}

The file contains only copyright notices.  It should either hold the
out-of-line definition of \code{AbsGame::kPass} (currently in the header)
or be deleted to reduce build noise.

\subsection{\style\ \ \file{lib/utils.cpp} — \code{qTrans()} is undocumented}

The function uses constants 3, 4, and 17 with no explanation of the
algorithm or its purpose (appears to be a fast non-cryptographic hash for
seed derivation).  A one-line comment stating the intent would prevent
future confusion.

\subsection{\style\ \ \file{src/Searcher.h}:16,21 — Two \code{mcts()} overloads with identical parameter types}

Both overloads take \code{const Game\&} and \code{int}; the second overload
adds another \code{int}.  A caller passing two ints unintentionally selects
the node-budget overload.  Strong typedef (\code{using Seconds = int}) or
distinct method names (\code{mctsByTime}, \code{mctsByNodes}) would prevent
accidental misuse.

\subsection{\style\ \ \file{lib/RectangularGraph.cpp} — Magic spacing constants}

\code{kColSpacing = 0.87f} and \code{kRowSpacing = 0.93f} are
unexplained.  A comment indicating these approximate
$\cos 30°$ (equilateral triangle spacing) and a slight vertical compression
would make the intent clear.

\subsection{\style\ \ \file{gui/BoardWidget.cpp}:216,256 — Repeated \code{dynamic\_cast} on every paint}

The same \code{dynamic\_cast<const IrregularGraph*>} is evaluated on every
\code{paintEvent()}.  The result could be cached in a \code{bool isIrregular\_}
member set once in \code{setGame()}, or replaced by a virtual method on
\code{Graph}.

\end{document}
% Copyright Ben Paul Wise. All Rights Reserved.
